%==================================================================
%  WHY THESE OPTIMIZERS BECAME FAMOUS
%  Beamer presentation — Madrid / whale theme
%  Matches existing Optimization & Linear Algebra deck style
%==================================================================

\documentclass{beamer}
\hypersetup{pdfpagemode=FullScreen}
\setbeamertemplate{navigation symbols}{}
\usepackage[english]{babel}
\usepackage[latin1]{inputenc}
\usepackage{amsmath,amssymb}
\usepackage[T1]{fontenc}
\usepackage{times}
\usepackage{mathptmx}
\usepackage{graphicx}
\usepackage{booktabs}
\usepackage{hyperref}
\usepackage{xcolor}
\usepackage{mathtools}

\mode<presentation>{
  \usetheme{Madrid}
  \usecolortheme{whale}
}

\setbeamertemplate{blocks}[rounded][shadow=false]
\AtBeginEnvironment{block}{\small}
\AtBeginEnvironment{alertblock}{\small}
\AtBeginEnvironment{exampleblock}{\small}
\setbeamersize{text margin left=6pt, text margin right=6pt}

\setbeamercovered{transparent}
\setbeamertemplate{bibliography item}[text]
\setbeamertemplate{theorems}[numbered]
\setbeamerfont{title}{size=\large}
\setbeamerfont{date}{size=\tiny}
\setbeamertemplate{itemize items}[ball]
\setbeamertemplate{itemize subitem}[circle]
\setbeamertemplate{itemize subsubitem}[triangle]
\setbeamertemplate{caption}[numbered]

%------------------------------------------------------------------
%  TITLE PAGE
%------------------------------------------------------------------
\title[Why Optimizers Became Famous]{%
  Why Did These Optimizers\\[4pt]
  Become Famous?}

\author[Vatsav Reddy]{%
  \texorpdfstring{%
    {\Large Vatsav Reddy}\\[6pt]
    {\small Under the Guidance of}\\[2pt]
    {\small Prof.\ Praveen Parchuri}%
  }{Vatsav Reddy}}

\institute[IIITH]{%
  {\normalsize International Institute of Information Technology, Hyderabad}}

\date{\today}

%==================================================================
\begin{document}
%==================================================================

\begin{frame}
  \titlepage
\end{frame}

\begin{frame}[shrink=10]{Contents}
  \tableofcontents
\end{frame}

\AtBeginSection[]{
  \begin{frame}<beamer>
    \frametitle{Current Section}
    \tableofcontents[currentsection]
  \end{frame}
}

%==================================================================
\section{The Question \& The Deep Answer}
%==================================================================

\begin{frame}[shrink=10]{Prof.\ Praveen's Question}
  \begin{alertblock}{The Question}
    Out of possibly \textbf{infinitely many} optimizers one could
    design mathematically, \textit{why did these specific few become
    famous?}
  \end{alertblock}

  \bigskip

  \begin{block}{The Surface Answer}
    \begin{itemize}\setlength\itemsep{5pt}
      \item \textbf{Not} because they have elegant equations
      \item \textbf{Not} because the papers were mathematically beautiful
      \item \textbf{Yes} because each solved a \emph{painful,
            real-world engineering bottleneck} in deep learning
    \end{itemize}
  \end{block}

  \bigskip

  \begin{exampleblock}{The Deep Scientific Answer}
    Almost all famous optimizers are approximating
    \textbf{Newton's method} cheaply under hardware, memory,
    and stochasticity constraints:
    \[
      w_{t+1} = w_t - H^{-1}\,\nabla\mathcal{L}(w_t)
    \]
    Storing $H$ costs $O(d^2)$; inverting costs $O(d^3)$~---
    \textbf{impossible} for deep learning.\\[4pt]
    The core question each optimizer answers is:\\
    \textit{``How can we imitate $H^{-1}\nabla\mathcal{L}$
    without ever computing $H$?''}
  \end{exampleblock}
\end{frame}

%------------------------------------------------------------------
\begin{frame}[shrink=10]{The Hidden Evolution of Optimizers}
  \begin{block}{Optimization Research as Failure-Mode Fixing}
    \textit{``What specific failure of the previous optimizer
    prevented neural networks from scaling?''}\\[6pt]
    Every famous optimizer fixed \textbf{one dominant failure mode}.
  \end{block}

  \medskip

  \begin{block}{The Evolutionary Chain}
    \centering
    \renewcommand{\arraystretch}{1.35}
    \begin{tabular}{lll}
      \toprule
      \textbf{Optimizer} & \textbf{Failure It Fixed} & \textbf{Intelligence Added}\\
      \midrule
      GD        & No optimizer existed              & Move downhill\\
      SGD       & Full-batch too expensive          & Compute at scale\\
      Momentum  & Zig-zagging, flat regions         & Remember the past\\
      NAG       & Blind overshooting                & Anticipate the future\\
      AdaGrad   & Sparse features under-trained     & Adapt per feature\\
      RMSProp   & Learning rate dies                & Forget old history\\
      Adam      & Too much manual tuning            & Combine all of the above\\
      AdamW     & Regularisation entangled in Adam  & Separate opt.\ from reg.\\
      \bottomrule
    \end{tabular}
  \end{block}
\end{frame}

%==================================================================
\section{Vanilla GD \& SGD: The Computation Bottleneck}
%==================================================================

\begin{frame}[shrink=10]{Vanilla GD --- Why It Became the Foundation}
  \begin{block}{Update Rule \& Hessian Interpretation}
    \[
      w_{t+1} = w_t - \eta\,\nabla\mathcal{L}(w_t)
              = w_t - \underbrace{\eta\,I}_{\approx\,H^{-1}}\,
                \nabla\mathcal{L}(w_t)
    \]
    GD assumes $H^{-1} \approx \eta I$~---
    \textbf{same curvature in every direction} (isotropic).
  \end{block}

  \medskip

  \begin{alertblock}{Why It Fails on Real Loss Surfaces}
    Neural network loss surfaces have highly \textbf{anisotropic}
    curvature: $\lambda_{\max} \gg \lambda_{\min}$
    (ill-conditioned Hessian).\\[4pt]
    Result: severe zig-zagging in narrow ravines;
    tiny gradients on gentle slopes $\Rightarrow$ very slow movement.
  \end{alertblock}

  \medskip

  \begin{block}{Why SGD Became Famous}
    \begin{itemize}\setlength\itemsep{5pt}
      \item Datasets became huge: ImageNet, web-scale NLP,
            recommendation systems
      \item Full-batch GD: \textbf{one update per epoch}
            (millions of calculations before any parameter move)
      \item SGD: \textbf{$N$ updates per epoch}~--- enabled
            online learning and memory-efficient training
      \item Made training on millions of samples \textbf{feasible}
            $\Rightarrow$ unlocked the deep learning boom
    \end{itemize}
  \end{block}

  \begin{exampleblock}{Real-World Bottleneck Solved}
    \textbf{Computational scalability}~--- not a mathematical
    improvement, a \emph{hardware and data reality} fix.
  \end{exampleblock}
\end{frame}

%==================================================================
\section{Momentum \& NAG: Curvature-Aware Dynamics}
%==================================================================

\begin{frame}[shrink=10]{Momentum --- Temporal Curvature Sensing}
  \begin{block}{Update Rule}
    \[
      u_t = \beta\,u_{t-1} + \eta\,\nabla w_t,
      \qquad
      w_{t+1} = w_t - u_t
    \]
    Expanding recursively:
    $u_t = \eta\displaystyle\sum_{i=1}^{t}\beta^{t-i}\,\nabla w_i$~---
    an \textbf{exponentially weighted moving average} of past gradients.
  \end{block}

  \medskip

  \begin{block}{Scientific Interpretation}
    \begin{itemize}\setlength\itemsep{5pt}
      \item \textbf{Flat valley (low curvature):} gradients stay
            \emph{aligned} for many steps
            $\Rightarrow$ EMA \emph{amplifies}~---
            optimizer accelerates
      \item \textbf{Steep direction (high curvature):} gradients
            \emph{flip sign} rapidly
            $\Rightarrow$ EMA \emph{cancels oscillations}
    \end{itemize}
    Momentum \textbf{implicitly senses curvature anisotropy}
    through temporal averaging~---
    a cheap approximation to second-order dynamics.
  \end{block}

  \begin{alertblock}{Real-World Bottleneck Solved}
    SGD zig-zagged in ravines;
    neural network loss surfaces are \textbf{not} spherical bowls~---
    they are ravines, curved valleys, anisotropic surfaces.
    Momentum matched this geometry.
    \textit{``Gain confidence, take bigger steps in same direction.''}
  \end{alertblock}
\end{frame}

%------------------------------------------------------------------
\begin{frame}[shrink=10]{Nesterov Accelerated Gradient --- Predictive Curvature Correction}
  \begin{columns}[T]
    \column{0.48\textwidth}
    \begin{block}{Momentum: Standing Here}
      \begin{enumerate}\setlength\itemsep{4pt}
        \item Stand at $w_t$
        \item Check slope \emph{here}
        \item Move (may overshoot)
      \end{enumerate}
    \end{block}

    \column{0.48\textwidth}
    \begin{block}{NAG: Look Ahead First}
      \begin{enumerate}\setlength\itemsep{4pt}
        \item Peek at $w_t - \beta u_{t-1}$
        \item Check \textbf{future} slope
        \item Adjust \emph{before} moving
      \end{enumerate}
    \end{block}
  \end{columns}

  \medskip

  \begin{block}{Scientific Interpretation: Hessian Appears Explicitly}
    Taylor-expand the look-ahead gradient:
    \[
      \nabla\mathcal{L}(w_t - \beta u_{t-1})
        \;\approx\;
        \nabla\mathcal{L}(w_t) - \beta\,H\,u_{t-1}
    \]
    The term $\beta H u_{t-1}$ is a \textbf{curvature correction}.
    NAG is thus a first-order approximation to
    second-order predictive correction~---
    this is the deep reason it outperforms plain Momentum.
  \end{block}

  \begin{exampleblock}{Real-World Bottleneck Solved}
    Deeper networks became unstable; momentum's blind acceleration
    caused many U-turns.
    NAG's lookahead gave \textbf{smarter, more stable acceleration}~---
    \textit{``Running fast while looking ahead.''}
  \end{exampleblock}
\end{frame}

%==================================================================
\section{AdaGrad \& RMSProp: Per-Feature Adaptivity}
%==================================================================

\begin{frame}[shrink=10]{AdaGrad --- Diagonal Hessian Approximation}
  \begin{block}{Motivation: Sparse Features}
    \[
      \nabla w^1 = \bigl(f(x)-y\bigr)\cdot f(x)\cdot(1-f(x))\cdot x^1
    \]
    If $x^1$ is sparse (mostly 0): $\nabla w^1 \approx 0$ for most
    inputs $\Rightarrow$ $w^1$ gets very few updates despite possibly
    being \emph{critical} (e.g.\ rare words: ``earthquake'', ``fraud'').
  \end{block}

  \medskip

  \begin{block}{AdaGrad Update \& Newton Connection}
    \[
      v_t = v_{t-1} + (\nabla w_t)^2,
      \qquad
      w_{t+1} = w_t
        - \underbrace{\frac{\eta}{\sqrt{v_t+\varepsilon}}}_{\approx\,D_t}
        \,\nabla w_t
      \;\approx\; w_t - \eta\,H^{-1}\nabla w_t
    \]
    where $D_t = \operatorname{diag}((\sqrt{v_t}+\varepsilon)^{-1})$
    is a \textbf{diagonal approximation to $H^{-1}$}.\\[4pt]
    For quadratic loss $\mathcal{L}=\tfrac{1}{2}w^THw$,
    large-curvature directions produce consistently large gradients,
    so $v_t$ estimates curvature magnitude~---
    \textbf{diagonal preconditioning from gradient statistics}.
  \end{block}

  \begin{alertblock}{Real-World Bottleneck Solved + Fatal Flaw}
    Solved sparse feature under-training at Google-scale NLP and
    ad-ranking.
    Fatal flaw: $v_t\!\uparrow\infty$
    $\Rightarrow$ effective $\eta\!\to 0$~---
    AdaGrad gets stuck near convergence.
  \end{alertblock}
\end{frame}

%------------------------------------------------------------------
\begin{frame}[shrink=10]{RMSProp --- Adaptive Local Curvature via EMA}
  \begin{block}{Fix: Exponential Moving Average of Squared Gradients}
    \[
      v_t = \beta\,v_{t-1} + (1-\beta)(\nabla w_t)^2
      \qquad (\beta > 0.9)
    \]
    \[
      \boxed{w_{t+1} = w_t
        - \frac{\eta}{\sqrt{v_t+\varepsilon}}\,\nabla w_t}
    \]
    Now $v_t \approx \mathbb{E}[g^2]$~--- the \textbf{local second
    moment} of gradients, not the unbounded cumulative sum.
  \end{block}

  \medskip

  \begin{block}{Fisher Information Connection}
    For many probabilistic models:
    \[
      \mathbb{E}[g^2] \;\approx\; \mathcal{F}
      \quad\text{(diagonal Fisher Information)}
    \]
    The Fisher Information is a \textbf{natural curvature measure} in
    probability space.
    RMSProp therefore approximates $H^{-1}$ through
    \emph{local curvature statistics}~---
    \textbf{does not kill the learning rate}.
  \end{block}

  \begin{exampleblock}{Real-World Bottleneck Solved}
    Neural networks are \textbf{non-stationary}: gradients from 1000
    steps ago should not dominate today.
    EMA decay introduces \emph{forgetting and recency awareness},
    matching real training dynamics.
    Practical success $>$ theoretical elegance.
  \end{exampleblock}
\end{frame}

%==================================================================
\section{Adam \& AdamW: The Default Standard}
%==================================================================

\begin{frame}[shrink=10]{Adam --- Stochastic Diagonal Quasi-Newton Method}
  \begin{block}{Adam $=$ RMSProp $+$ Momentum $+$ Bias Correction}
    \begin{align*}
      m_t &= \beta_1\,m_{t-1} + (1-\beta_1)\,\nabla w_t
            &\leftarrow&\;\text{1st moment (direction)}\\
      v_t &= \beta_2\,v_{t-1} + (1-\beta_2)\,(\nabla w_t)^2
            &\leftarrow&\;\text{2nd moment (curvature)}\\[6pt]
      \hat{m}_t &= \frac{m_t}{1-\beta_1^t},
      \quad\hat{v}_t = \frac{v_t}{1-\beta_2^t}
            &\leftarrow&\;\text{bias correction}\\[4pt]
      \beta_1 &= 0.9,\quad\beta_2 = 0.999,\quad\varepsilon=10^{-8}
    \end{align*}
    \[
      \boxed{w_{t+1} = w_t
        - \frac{\eta}{\sqrt{\hat{v}_t}+\varepsilon}\cdot\hat{m}_t}
      \;=\; w_t - \eta\,\underbrace{D_t^{-1}}_{\approx H^{-1}}\,\hat{m}_t
    \]
    \textbf{Adam is a stochastic, diagonal, quasi-Newton method}
    with momentum.
  \end{block}

  \begin{alertblock}{Real-World Bottleneck Solved}
    Before Adam: careful per-task LR tuning, optimizer instability.
    After Adam: one set of defaults works across
    CNNs, RNNs, Transformers, GANs, and diffusion models.
    \textbf{Engineering practicality beats theoretical purity}~---
    even though Adam does not converge in some cases, industry
    adopted it universally because it \emph{reduces human suffering}.
  \end{alertblock}
\end{frame}

%------------------------------------------------------------------
\begin{frame}[shrink=10]{AdamW --- Decoupled Geometry and Regularisation}
  \begin{block}{The Hidden Problem in Adam + L2}
    Standard Adam with L2 regularisation sets
    $g_t \leftarrow \nabla\mathcal{L}(w_t) + \lambda w_t$.
    Now $\lambda w_t$ enters the curvature accumulator $v_t$,
    \textbf{entangling regularisation with adaptive scaling}~---
    weights that are large get a smaller effective learning rate,
    which is not true weight decay.
  \end{block}

  \medskip

  \begin{block}{AdamW Fix: Apply Weight Decay After the Adaptive Step}
    \[
      w_{t+1}
        = w_t
        - \frac{\eta\,\hat{m}_t}{\sqrt{\hat{v}_t}+\varepsilon}
        - \underbrace{\eta\lambda\,w_t}_{\text{true weight decay}}
    \]
    \begin{itemize}\setlength\itemsep{5pt}
      \item \textbf{Optimisation geometry}: $\hat{m}_t/\sqrt{\hat{v}_t}$
      \item \textbf{Regularisation geometry}: $\lambda w_t$
            (independent of curvature scaling)
    \end{itemize}
    Cleaner decoupling $\Rightarrow$ better generalisation and
    stability for over-parameterised models.
  \end{block}

  \begin{exampleblock}{Real-World Bottleneck Solved}
    Transformers are hugely over-parameterised and
    sensitive to regularisation.
    AdamW became \textbf{the standard for LLMs and ViTs} because
    decoupled weight decay improved generalisation at scale.
  \end{exampleblock}
\end{frame}

%==================================================================
\section{Grand Unification: Approximating Newton Cheaply}
%==================================================================

\begin{frame}[shrink=10]{Grand Unification: Every Optimizer Approximates $H^{-1}$}
  \begin{block}{Newton's Ideal (Unachievable in Practice)}
    \[
      w_{t+1} = w_t - H^{-1}\nabla\mathcal{L}(w_t)
      \qquad
      \bigl[\text{cost: }O(d^2)\text{ storage},\;
             O(d^3)\text{ inversion}\bigr]
    \]
  \end{block}

  \medskip

  \begin{block}{The $O(d)$ Approximation Hierarchy}
    \centering
    \renewcommand{\arraystretch}{1.35}
    \begin{tabular}{lll}
      \toprule
      \textbf{Optimizer} & \textbf{$H^{-1}$ Approximation} & \textbf{Key Insight}\\
      \midrule
      GD       & $\eta I$                              & Isotropic; ignores curvature\\
      Momentum & EMA filter on gradients               & Temporal curvature sensing\\
      NAG      & $I - \beta H$ (1st-order)             & Curvature correction via lookahead\\
      AdaGrad  & $\operatorname{diag}(\textstyle\sum g^2)^{-1/2}$
                                                       & Cumulative diagonal preconditioning\\
      RMSProp  & $\operatorname{diag}(\mathbb{E}[g^2])^{-1/2}$
                                                       & Local curvature $\approx$ Fisher info\\
      Adam     & EMA diagonal + smoothed gradient      & Quasi-Newton with momentum\\
      AdamW    & Adam + decoupled $\lambda w$          & Separate geometry \& regularisation\\
      \bottomrule
    \end{tabular}
  \end{block}
\end{frame}

%------------------------------------------------------------------
\begin{frame}[shrink=10]{Why Only These Survive: The Brutal Filter}
  \begin{block}{Survival Criteria~--- All Must Be Satisfied}
    \begin{itemize}\setlength\itemsep{6pt}
      \item \textbf{$O(d)$ cost}: scales to billions of parameters
            on GPU hardware
      \item \textbf{Stochastic robustness}: stable under mini-batch
            gradient noise, no full-batch requirement
      \item \textbf{Minimal tuning}: default hyperparameters
            transfer across tasks without per-problem search
      \item \textbf{GPU parallelism}: element-wise operations;
            no sequential or high-memory-bandwidth steps
      \item \textbf{Predictable behaviour}: converges (or diverges
            predictably) on standard architectures
      \item \textbf{Reduces human effort}: lowers the number of
            manual engineering decisions during training
    \end{itemize}
  \end{block}

  \begin{alertblock}{The Core Insight}
    There are \emph{infinitely many} mathematically valid optimizers.
    Most fail at least one criterion above.
    The few that satisfy \textbf{all criteria simultaneously}
    are the names you know.\\[4pt]
    \textbf{The history of optimizers is the history of cheap
    approximations to $H^{-1}$ under real-world constraints.}
  \end{alertblock}
\end{frame}

%==================================================================
\section{Summary}
%==================================================================

\begin{frame}[shrink=10]{Summary}
  \begin{columns}[T]
    \column{0.50\textwidth}
    \begin{block}{Engineering View}
      \begin{itemize}\setlength\itemsep{4pt}
        \item \textbf{GD / SGD}: feasibility at data scale;
              $N$ updates per epoch vs.\ $1$
        \item \textbf{Momentum}: smooth traversal of anisotropic
              loss surfaces; EMA velocity
        \item \textbf{NAG}: lookahead curvature correction;
              reduces U-turns in deep nets
        \item \textbf{AdaGrad}: per-feature adaptive $\eta$;
              sparse feature parity
        \item \textbf{RMSProp}: non-stationary training dynamics;
              EMA prevents $\eta\!\to 0$
        \item \textbf{Adam}: plug-and-play default; combines
              momentum + adaptivity + bias correction
        \item \textbf{AdamW}: Transformer regularisation;
              decoupled weight decay
      \end{itemize}
    \end{block}

    \column{0.46\textwidth}
    \begin{block}{Scientific View ($H^{-1}$ Approx.)}
      \begin{itemize}\setlength\itemsep{4pt}
        \item \textbf{GD}: $H^{-1}\approx\eta I$
        \item \textbf{Momentum}: temporal curvature
              sensing via gradient EMA
        \item \textbf{NAG}: $\nabla\mathcal{L}(w-\beta u)$
              $\Rightarrow Hu$ term appears
        \item \textbf{AdaGrad}: cumulative diagonal
              $H^{-1}$ estimator
        \item \textbf{RMSProp}: EMA diagonal
              $\approx$ Fisher information
        \item \textbf{Adam}: stochastic diagonal
              quasi-Newton
        \item \textbf{AdamW}: Adam with decoupled
              curvature \& regularisation
      \end{itemize}
    \end{block}
  \end{columns}

  \medskip

  \begin{exampleblock}{One-Line Answer for Prof.\ Praveen}
    These optimizers became famous \textbf{not} because they were
    mathematically possible, but because each solved a
    \textbf{dominant real-world bottleneck}~--- computation,
    instability, sparse features, convergence speed, adaptivity,
    or generalisation~--- while remaining $O(d)$ and robust enough
    to use as a \textbf{default}.
  \end{exampleblock}
\end{frame}

%------------------------------------------------------------------
\begin{frame}
  \Huge{\centerline{Thank You}}
  \vspace{0.6cm}
  \large{\centerline{Questions?}}
\end{frame}

%==================================================================
\end{document}
%==================================================================